\documentclass[conference]{IEEEtran}
\IEEEoverridecommandlockouts
\usepackage{cite}
\usepackage{amsmath,amssymb,amsfonts}
\usepackage{graphicx}
\usepackage{textcomp}
\usepackage{xcolor}
\usepackage{url}
\usepackage[hidelinks]{hyperref}
\renewcommand{\IEEEkeywordsname}{Keywords}
\emergencystretch=3em
\usepackage{array}
\begin{document}

\title{Fleet on the Floor: A Fog-to-Cloud Pipeline for Robotics Health Monitoring}

\author{\IEEEauthorblockN{Goutham Uppu}
\IEEEauthorblockA{Student ID: X25167936 \\ MSc in Cloud Computing \\ Fog and Edge Computing (H9FECC) \\ National College of Ireland \\ x25167936@student.ncirl.ie}}

\maketitle
\bstctlcite{BSTcontrol}

\begin{abstract}
Warehouse robots move continuously across the floor, yet their health is often known only from nominal schedules rather than live condition data. This paper presents a continuous fleet-monitoring pipeline for ten autonomous mobile robots across two warehouse zones. Five telemetry channels per robot feed a fog gateway that windows and aggregates readings and evaluates threshold checks locally, forwarding only compact summaries to the cloud. A scalable serverless backend on Amazon SQS, AWS Lambda, and DynamoDB ingests those summaries and serves a live dashboard through Amazon API Gateway, with the static frontend hosted on Amazon S3 and the whole cloud tier provisioned in a single Terraform apply. Each layer is justified against alternatives, and the system is deployed to a real AWS account rather than a local emulator. An automated suite of 116 tests passes, and end-to-end verification against the live deployment confirmed every pipeline health check reporting healthy and live fleet data rendering correctly in a real browser. A pre-deployment audit additionally identified and corrected three defects, in pagination, message batching, and credential handling, that the test suite alone had not surfaced.
\end{abstract}

\begin{IEEEkeywords}
fog computing, autonomous mobile robots, warehouse automation, Amazon SQS, AWS Lambda, DynamoDB, Infrastructure as Code, Terraform, fleet health monitoring
\end{IEEEkeywords}


\section{Introduction}

Autonomous mobile robots (AMRs) have moved from pilot projects to a standard fixture on warehouse floors, ferrying totes and pallets between racking, picking stations and dock doors without a fixed rail to follow. A shop-floor fleet-management system built for exactly this setting found that dispatch decisions improve once a robot's position, battery state and current task load are known live, not just its nominal schedule~\cite{ref1}. A recent survey of intralogistics AMR deployments reaches a related conclusion from the opposite direction: it names resource scheduling, power consumption and safety as gaps the field has not closed, in large part because so few deployments instrument the robots closely enough to reason about all three together~\cite{ref2}. Warehouse Robotics Fleet Monitoring targets that gap directly for a fleet of ten AMRs split across two zones, zone-a and zone-b, giving each robot five independent telemetry channels (battery level, payload weight, motor temperature, position drift, and task-queue depth) and a pipeline that turns those channels into a live operational picture instead of a log read after a breakdown.

Bonomi et al. define fog computing as an extension of the cloud paradigm out to the edge of the network, placing storage, compute and networking on nodes that sit physically close to where data originates~\cite{ref3}. That definition describes this project's own topology closely: a fog gateway runs on the same local segment as the ten sensor containers, windowing their readings and evaluating fleet-health thresholds before anything leaves the building. The cloud tier that eventually receives the aggregated result needs its own definition, and the NIST reference model supplies one: on-demand, broadly accessible, elastically provisioned services metered by use~\cite{ref4}. DynamoDB, SQS, Lambda and API Gateway all satisfy that description, and the boundary between the two tiers is treated as a design decision: what crosses it is a windowed aggregate, not a raw sample, because the fog node exists specifically to make that reduction before anything crosses the network.

The project's objective was to build and deploy a complete fog-to-cloud pipeline for this fleet, not to simulate one against a local emulator and stop there. Ten sensor containers, five sensor types across two zones, generate readings on independent sample and dispatch intervals; a fog gateway ingests, buffers and windows them per robot and per zone, evaluates a set of alert rules against configurable thresholds, and relays the aggregates onward in batches; a queue decouples the fog tier from the cloud tier so that neither side blocks on the other; a Lambda function consumes the queue into a DynamoDB table; and a dashboard, backed by a second Lambda function behind a real API Gateway, renders a fleet-ops view for whoever is watching the floor.

Meeting that objective set a specific bar: a believable multi-sensor simulation with bounded, physically plausible drift instead of uniform random noise; a fog tier that does real aggregation work (windowing, threshold evaluation, batching) instead of forwarding raw samples unchanged; a scalable cloud tier built from managed AWS services rather than a single long-running server; a dashboard an operator could actually read at a glance; automated tests covering the arithmetic and edge cases that matter, not just the happy path; a continuous-integration pipeline that runs those tests on every push rather than trusting a developer to run them locally; and evidence the system runs somewhere real, an actual deployment to AWS rather than a description of one.

The remainder of this report is organised as follows. Section II sets out the architecture and the reasoning behind its more contested choices. Section III describes the software as implemented, the Terraform module that provisions the entire cloud tier in one operation, the results of the automated test suite, and the real AWS deployment, including the deployment environment, the defects identified and corrected before deployment, and the evidence gathered against the running system. Section IV concludes with an interpretation of those findings and a short reflection on building the project.

\section{Architecture}

Each of the ten sensor containers represents one robot in one zone and publishes its five metrics on two independently configurable rates, one controlling how often a new reading is generated and the other how often a batch of readings is pushed to the fog gateway, so that noisy, frequent sampling can be decoupled from network chatter. The fog gateway is a plain Java HTTP server instead of a framework-managed service: it accepts ingest posts, buffers them per robot and per zone, and closes a window on a fixed interval, at which point it computes an aggregate and checks it against a set of alert rules. Battery telemetry receives particular attention in that rule set, because predictive battery monitoring for AMR fleets has been shown to catch degradation trends well before an outright failure, using continuous IIoT-style telemetry of exactly the kind this project collects instead of a periodic manual check~\cite{ref5}. The alert rules for motor temperature and position drift follow the same pattern: a threshold crossing is flagged at the window level, where the dashboard's detail panel can show the trend rather than an operator having to scroll a per-sample log.

On the cloud side, DynamoDB's own partition-key guidance sets a hard limit of 3,000 read-capacity units and 1,000 write-capacity units per partition, and recommends spreading writes so no single partition absorbs a disproportionate share of the fleet's traffic~\cite{ref6}. The table's key design, a sort key disambiguated by the window's end timestamp and the site identifier, follows that guidance by keying on the sensor/robot dimension instead of a single fleet-wide counter, so ten robots' windows land across the partition space rather than piling onto one key. SQS sits between the fog gateway and the ingestion function for the usual reason a queue separates a producer from a consumer that should not share a failure mode, but the batching decision built on top of it has a specific justification: SQS bills and rate-limits per request, and its batch API amortises that cost across up to ten messages instead of one request per message~\cite{ref7}. The relay publisher chunks a window's aggregates at exactly that ten-entry limit, called once per closed window instead of once per aggregate, the difference this project's own load test, described in Section III, was built to demonstrate.

The overall pipeline shape (fog to queue to Lambda to DynamoDB, queried by a second Lambda behind API Gateway), drawn end to end in Fig.~\ref{fig:topology}, is a widely used serverless pattern; what is distinctive here is how the dashboard function models its routes. Each route is a distinct variant of a single sealed Java type, so the compiler can verify that every possible route is handled: adding a new route without handling it is a compile-time error rather than a not-found response discovered at runtime, a stronger guarantee than a lookup table or a reflection-driven router provides. The cost is a dispatch that is extended by hand for every new route instead of one populated declaratively, but the API surface here is five fixed routes unlikely to grow substantially, so the compile-time guarantee mattered more than declarative extensibility would have for a surface this size.

\begin{figure*}[!t]
\centering
\includegraphics[width=.9\textwidth]{topology.png}
\caption{Deployment topology. Sensors and Fog Gateway sit inside the EC2 warehouse-aisle loop's top-left corner; the ingest path runs clockwise through SQS, wrf-processor and DynamoDB, and the query path runs counter-clockwise back through wrf-dashboard-api, API Gateway/S3 and the Browser. Three amber dotted spokes mark wrf-dashboard-api's health/state reads (SQS queue depth, Fog Gateway, wrf-processor); the fourth, from DynamoDB, is drawn as a solid arrow since it doubles as the ingest path's own final edge.}
\label{fig:topology}
\end{figure*}

wrf-dashboard-api's health handler reasons about pipeline freshness by reading four things directly instead of trusting a single cached flag: whether the fog gateway's own health endpoint answers over HTTP, whether the SQS queue is reachable, whether the wrf-processor Lambda's own deployment state reports Active, and how old the most recent DynamoDB window is across every sensor type. That last check, the freshest-window age, is the one that actually detects a stalled pipeline: a queue and a Lambda can both report healthy while no new data has landed in over an hour, and only a timestamp comparison against the table itself catches that. The fog gateway's health and thresholds endpoints answer without any credential check, which reflects a scope decision specific to this deployment: they return only aggregate window state and configured threshold values, never a live per-robot reading, so what they expose is bounded to information an operator already sees rendered on the dashboard itself.

Two alternatives were weighed and set aside. Kinesis Data Streams would have kept an ordered, replayable record of every window per shard, which SQS's standard queue does not guarantee, but ordering was judged unnecessary here: the dashboard reads the latest window per sensor type from DynamoDB directly, not a stream of intermediate states, so Kinesis's stronger ordering guarantee would have added shard-management overhead without a matching need to spend it on. A single Lambda function handling both ingestion and dashboard queries was also considered and set aside, because ingestion is triggered by SQS on the fleet's own schedule while dashboard queries are triggered by whoever is looking at the browser: coupling the two would mean a slow dashboard query competing for the same concurrency budget as message processing, and splitting them into separate ingestion and dashboard functions keeps those two very different load profiles from interfering with each other.

\section{Implementation}

The entire pipeline is written in Java 17 against the plain Java HTTP server, with no web framework in any of the four modules: sensors, fog, processor, and dashboard. Each module builds and tests independently; the dashboard module additionally produces a second, Lambda-facing entry point alongside its local development server, so the same query and threshold logic serves both a laptop and a real API Gateway integration.

A GitHub Actions workflow runs on every push and pull request as two dependent jobs rather than one. The first, unit-tests, runs the Maven test goal in each of the four modules in turn (sensors, fog, backend/processor and backend/dashboard), so a failure in any one module's own test suite fails the build before the second job starts. The second, integration, brings the full stack up with Docker Compose against LocalStack and runs the pipeline verification script against the running containers, exercising the same ingest-to-DynamoDB path this report verifies separately against real AWS later in this section, then tears the stack back down regardless of outcome so a failed run does not leave containers running on the runner.

Each sensor process simulates one robot's one sensor type, using a bounded random walk to keep each metric's drift in range instead of letting it wander arbitrarily between samples: position drift, for instance, is walked around a centre value rather than accumulating without limit, which matters because an unbounded drift value would eventually trip every threshold rule regardless of whether the robot is actually behaving abnormally. Ten containers run in total, one per zone/sensor-type combination across zone-a and zone-b, each independently configurable on its own sample and dispatch intervals. Each dispatch posts a small JSON object naming the sensor type, zone identifier and unit alongside the batch of timestamped values collected since the previous send; at the default two-second sample and ten-second dispatch intervals each batch carries about five readings, and the whole payload comes to a few hundred bytes, well under a kilobyte.

The fog gateway ingests posted readings over HTTP, buffers them per robot and per zone so that one zone's burst of traffic cannot delay another zone's window from closing, computes a window aggregate when the window elapses, evaluates that aggregate against the configured alert rules, and calls the relay publisher to send the result onward. If that outbound send to the queue fails, the window's aggregates, already drained from the per-robot buffers by that point, are not retried or held anywhere else: the fog gateway logs the failure and moves on to the next window, an accepted loss at this data volume given how soon the next window closes. A separate thresholds endpoint exposes the currently configured rule values so the dashboard can display them next to the readings they gate, instead of hardcoding a second copy of the same numbers into the frontend.

The processor's Lambda entry point is invoked by an SQS event-source mapping; a pure transform step builds the table's composite sort key and maps a batch of queue records into DynamoDB writes, with per-record failures tallied and returned as a partial batch-item failure response instead of one failure discarding the whole batch.

Every AWS resource this project runs on was provisioned by a single Terraform module, applied once against Goutham Uppu's own Learner Lab account. The module's 24 resources span every service group the deployment needs: the DynamoDB table; the SQS queue; both Lambda functions plus the event source mapping connecting the queue to wrf-processor; the full API Gateway REST API chain (a resource, a proxy method and integration at the root and at the greedy proxy path, a Lambda invoke permission, a deployment and a prod stage); the EC2 instance, its security group and its Elastic IP; and the two S3 buckets, the frontend bucket's public-access block, bucket policy, website configuration and object upload, plus a provisioning step that ships the fog/sensor source tarball into the EC2 instance's user-data script on first boot. An empirical study of infrastructure-as-code adoption across 44 companies found that the practice's main payoff is reproducibility: the same declared state can be torn down and rebuilt without hand-run steps drifting between environments~\cite{ref8}, and this project realises that payoff directly: rather than a sequence of individual AWS CLI commands run in order, one Terraform apply brought all 24 resources up with matching configuration, down to the dashboard Lambda's fog-health and fog-thresholds environment variables being wired automatically to the Elastic IP Terraform itself allocated, with no manual copying of an IP address between commands. That convenience has a cost worth naming: the module couples the frontend upload and the EC2 bootstrap to the same apply, so a change touching only the dashboard's static assets still triggers a plan against the entire stack, and a later iteration would benefit from splitting the frontend deploy into its own narrower Terraform workspace.

The project's source is maintained at \href{[GitHub repository URL]}{[GitHub repository URL]}.

The automated suite comprises 116 tests across the four Maven modules, summarised by module and focus in Table~\ref{tab:tests}; all pass. Coverage emphasises bounded-drift behaviour in the sensors, exactly-at-boundary threshold evaluation in the fog tier, partial-batch-failure handling in the processor, and the multi-page Scan-pagination path in the dashboard. No test in the suite touches a real AWS account or a local emulator; each AWS client interface is backed by a hand-written fake.

\begin{table}[!t]
\caption{Automated Test Suite by Module (116 Tests Total)}
\label{tab:tests}
\centering
\footnotesize
\begin{tabular}{|>{\raggedright\arraybackslash}p{1.75cm}|p{0.7cm}|>{\raggedright\arraybackslash}p{4.45cm}|}
\hline
\textbf{Module} & \textbf{Tests} & \textbf{Focus} \\
\hline
sensors & 56 & bounded random walk (53), dual-rate sample/dispatch loops (3) \\
\hline
fog & 24 & buffer, alert rules, exactly-at-boundary threshold checks \\
\hline
backend/processor & 9 & sort-key mapping, partial-batch-failure tallying \\
\hline
backend/dashboard & 27 & routes, Scan-pagination fix, Lambda dispatch tests \\
\hline
\end{tabular}
\end{table}

Scalability evidence beyond the unit level comes from a burst-test script, which drives 2,000 messages through 32 concurrent workers against the running stack. Its purpose is narrower than a full soak test: it demonstrates that the fog-to-queue path holds up under concurrent load from many simulated robots at once, independent of whatever cadence the ten containers themselves are configured to run at.

The deployed dashboard was checked two ways. curl against the health, backend-stats, fleet and thresholds endpoints confirmed the API Gateway integration answers correctly, and successive backend-stats polls showed the DynamoDB item count increasing over time, consistent with sensor data actively arriving instead of a table that had stopped updating. curl alone cannot confirm that a browser loading the dashboard's S3 URL can actually reach that API, parse the response and render it: a CORS misconfiguration or a broken static-asset path returns a perfectly valid response to curl while leaving the dashboard blank in every browser that loads it, a failure mode that is easy to miss precisely because every API-level check passes. A Playwright session against the live URL confirmed the fleet-ops HUD (the roster table with its inline sparklines and LED status indicators, and the detail panel beneath it) renders real, changing sensor data with no console error beyond a single missing-favicon 404, which has no bearing on the pipeline. That view refreshes on a fixed polling loop rather than a server push: the frontend re-fetches the fleet, health and backend-stats endpoints together every two and a half seconds and re-renders the roster, sparklines, detail panel and footer from whatever those three responses currently hold. Fig.~\ref{fig:dashboard} reproduces the view that session captured against the live deployment.

\begin{figure}[!t]
\centering
\includegraphics[width=\columnwidth]{dashboard.png}
\caption{The deployed fleet-ops dashboard at its S3 URL, backed by wrf-dashboard-api through API Gateway: the two-zone robot roster with inline sparklines and LED status indicators, the per-robot telemetry detail panel, and all four pipeline health checks reporting healthy.}
\label{fig:dashboard}
\end{figure}

Warehouse Robotics Fleet Monitoring runs in Goutham Uppu's own AWS Academy Learner Lab account, region us-east-1, under the account's session role, reusing the shared lab role and instance profile instead of any role created specifically for this project: the account cannot create new IAM roles or users, a standard constraint of AWS Academy Learner Lab accounts.

A pre-deployment code audit examined the item-count method the backend-stats endpoint calls to report how many items sit in the readings table, and found it issued a single scan and returned that scan's count directly: correct only while the table holds fewer items than one scan page returns, and wrong once it grows past that without any indication in the returned number. The fix walks every page of the scan, following the pagination cursor from one page to the next and summing the per-page counts, so the reported total stays correct regardless of table size.

The same audit found the fog gateway relaying each window's aggregates to SQS one send at a time. The relay publisher now chunks a window's payloads into groups of up to ten and issues one batch send per chunk, called once per closed window instead of once per aggregate.

The most consequential defect was in credential handling, present in all three AWS-facing classes: the processor's client builder, the local dashboard server's client builder, and the relay publisher's constructor all built a static test-credentials provider unconditionally, instead of only when an explicit LocalStack endpoint was configured. In the two client builders that would have meant a real Lambda invocation quietly overriding its own execution-role credentials with a pair AWS would reject outright. In the relay publisher the same bug was worse: the endpoint override was also applied unconditionally in the original code, and the endpoint value is null outside a LocalStack environment, so the fog node would have thrown a null-pointer exception and failed to start at all the first time it ran against real AWS, not merely authenticated incorrectly. All three were fixed to gate on the endpoint variable's presence, matching the pattern the codebase already used correctly in one place.

\begin{table}[!t]
\caption{Live AWS Resources (us-east-1)}
\label{tab:resources}
\centering
\footnotesize
\begin{tabular}{|>{\raggedright\arraybackslash}p{1.7cm}|>{\raggedright\arraybackslash}p{2.9cm}|>{\raggedright\arraybackslash}p{2.6cm}|}
\hline
\textbf{Resource} & \textbf{Identifier} & \textbf{Role} \\
\hline
DynamoDB table & wrf-readings & stores closed-window aggregates \\
\hline
SQS queue & wrf-fleet-agg & decouples fog gateway from wrf-processor \\
\hline
Lambda & wrf-processor & SQS-triggered ingestion into DynamoDB \\
\hline
Lambda & wrf-dashboard-api & sealed-type route dispatch behind API Gateway \\
\hline
API Gateway REST API & iodllqqk3m & public entry point for the dashboard API \\
\hline
EC2 instance & i-00c6537b8a41e9750 & fog gateway + 10 sensor containers \\
\hline
Elastic IP & 3.211.126.248 & fixed address for the EC2 instance \\
\hline
S3 buckets & wrf-frontend-789399341650 / wrf-deploy-789399341650 & static dashboard frontend / deploy staging \\
\hline
\end{tabular}
\end{table}

With those three defects corrected, the stack was provisioned in a single Terraform apply against the confirmed account (24 resources, described earlier in this section, with zero manual AWS CLI steps in between). The full set of provisioned resources, with identifiers and roles, is listed in Table~\ref{tab:resources}. Verification went beyond the application's own self-report: the health endpoint was polled directly and returned all four fields true, backend-stats showed the DynamoDB item count climbing across successive polls, and the dashboard was opened in a Playwright-driven browser session at its S3 URL, where it rendered the fleet roster and detail panel with live data and no console error beyond the one harmless missing-favicon 404 already noted earlier in this section.

\section{Conclusion}

The project set out to build and deploy a complete fog-to-cloud pipeline for a ten-robot warehouse fleet, and each layer of that objective was met: a dual-rate multi-sensor simulation with bounded drift, a fog tier that windows readings, evaluates thresholds, and batches its output rather than relaying raw samples, a scalable serverless backend of managed AWS services, and a deployment verified end to end against a real AWS account and a real browser rather than a local emulation.

All three defects the pre-deployment audit turned up in Warehouse Robotics Fleet Monitoring shared the same shape: a Scan response's per-page count versus the whole table's count, a single SQS publish versus a batched one, and a credential provider built for LocalStack versus one built for a real Lambda execution role. None of the 116 unit tests that passed before the audit began would have caught any of the three, because none of them exercised more than one Scan page, more than one publish call, or a real endpoint-absent code path. That gap is the practical argument for treating a code audit against real service semantics as a distinct step from running an existing test suite, not a substitute for one.

The Terraform module is the other substantive outcome: a CLI-scripted deployment assembles a cloud tier by running AWS commands in sequence, copying an Elastic IP's value from one command's output into the next command's environment variable by hand. Collapsing that into one declared module and one apply removes the specific failure mode where a manually copied IP address is stale or mistyped, and gives the deployment a single state file recording what was actually created.

A key lesson concerns the limits of a passing test suite: the pagination and credential defects each passed every existing test, because those tests shared the assumptions that produced the bugs. This argues for validating against real service semantics rather than relying on a local emulator alone. Ordering dependencies within the Terraform module, in particular the EC2 user-data script's reliance on the S3 upload completing first, were the most demanding aspect to provision correctly, and the session-limited deployment account constrained how much could be verified per iteration.

One item remains open for future work: the Terraform module's coupling between the frontend upload and the EC2 bootstrap, noted in Section III, would be worth splitting once the deployment needs to change more often than it does today. Everything reported here reflects what the code, the test suite, and the live AWS account show today.

\bibliographystyle{IEEEtran}
\bibliography{refs}

\end{document}
